% ============================================================
%  LECSEG — 8-page IEEE-style conference paper
%  Use IEEEtran.cls (place IEEEtran.cls + IEEEtran.bst in this folder
%  if not installed system-wide).
% ============================================================

\documentclass[conference]{IEEEtran}
\IEEEoverridecommandlockouts

\usepackage{cite}
\usepackage{amsmath, amssymb, amsfonts}
\usepackage{algorithmic}
\usepackage{graphicx}
\usepackage{textcomp}
\usepackage{xcolor}
\usepackage{booktabs}
\usepackage[hidelinks]{hyperref}

\def\BibTeX{{\rm B\kern-.05em{\sc i\kern-.025em b}\kern-.08em
    T\kern-.1667em\lower.7ex\hbox{E}\kern-.125emX}}

\newcommand{\lecseg}{LecSeg}
\newcommand{\dataset}{LecSeg-30}
\newcommand{\todo}[1]{\textcolor{red}{\textbf{[#1]}}}

\begin{document}

\title{Hierarchical Multimodal Lecture-Video Topic\\Segmentation with
       Reliability-Weighted Fusion}

\author{
  \IEEEauthorblockN{Author One, Author Two, Author Three, Author Four, Author Five}
  \IEEEauthorblockA{Department of Computer Science and Engineering\\
    BRAC University, Dhaka, Bangladesh\\
    Email: \{firstname.lastname\}@bracu.ac.bd}
}

\maketitle

\begin{abstract}
We present \lecseg{}, an open hierarchical multimodal pipeline for
lecture-video topic segmentation that produces both chapter-level and
subtopic-level boundaries. \lecseg{} introduces (i) a reliability-weighted
fusion module that adapts per video and per sentence to noisy modalities, (ii)
a joint two-level hierarchy head, and (iii) a boundary refinement and titling
stage based on a small open local language model. We release \dataset{}: a
30-video, 5-domain corpus with dual subtopic annotation. Across five
established metrics with 1{,}000-sample bootstrap CIs and paired Wilcoxon
testing, \lecseg{} reduces $P_k$ by \todo{X.XX}\,\% relative over the
strongest open baseline ($p<\todo{0.0X}$). All code, dataset metadata, and
trained weights are publicly released.
\end{abstract}

\begin{IEEEkeywords}
topic segmentation, lecture video, multimodal fusion, hierarchical
segmentation, local language models.
\end{IEEEkeywords}

% ----------
\section{Introduction}
\todo{1 column. Motivation, gap, contributions (4 bullets), paper structure.}

% ----------
\section{Related Work}
\todo{1 column. 4 paragraphs: classical, neural, multimodal, hierarchical.}

% ----------
\section{The \lecseg{} Pipeline}
\todo{1.5 columns + Fig.~1. Stage diagram, fusion module, hierarchy head,
LLM refinement.}

% ----------
\section{The \dataset{} Corpus}
\todo{0.5 column. Source criteria, annotation protocol, IAA.}

% ----------
\section{Experiments}
\todo{1.5 columns + Table I (main) + Table II (ablations). Metrics,
significance tests, error categories.}

\begin{table}[t]
  \centering
  \caption{Main results on \dataset{} (mean $\pm$ 95\,\% bootstrap CI).}
  \label{tab:ieee_main}
  \begin{tabular}{lccc}
    \toprule
    Method & $P_k\downarrow$ & WD$\downarrow$ & BS$\uparrow$ \\
    \midrule
    TextTiling   & \todo{} & \todo{} & \todo{} \\
    BERT-Seg     & \todo{} & \todo{} & \todo{} \\
    \lecseg{} (ours) & \textbf{\todo{}} & \textbf{\todo{}} & \textbf{\todo{}} \\
    \bottomrule
  \end{tabular}
\end{table}

% ----------
\section{Conclusion}
\todo{4 sentences. Summary + 1 future direction.}

% ----------
\section*{Acknowledgements}
\todo{Supervisor + dept.}

\bibliographystyle{IEEEtran}
\bibliography{../thesis/bibliography/references}

\end{document}
